\section{Deferred Proofs}
\label{sec:app-proofs}

This appendix records commentary on the proofs whose main ideas already appear
in the body.

\subsection{Uniform perturbation calculus}

The proof of \Cref{prop:4.1} is intentionally elementary.  Its value is not
technical depth but category control: replace-one is an operation composition,
whereas LOO is an operation-plus-evaluation convention.  The triangle-inequality
argument therefore applies directly to delete/add versus replace, but not to LOO
without additional translation lemmas.

\subsection{Margin-crossing criterion}

The proof of \Cref{prop:5.1} reduces the classifier to a signed vote threshold.
All geometric complexity is compressed into the ordered top-$k$ neighbor list
$N_k(S,x)$.  Once that list is fixed, the classifier is just the sign of
$\M_k(S,x)$ with the non-positive side resolved toward label $0$.  The only way
a perturbation changes the prediction is therefore to move the signed margin
across zero.

\subsection{Two-point $1$-NN witness}

The witness of \Cref{sec:minimal-1nn} is fully explicit and hand-checkable.  No
exhaustive search is needed for the existence claim itself.  What computational
search contributes is different: it shows that within the enumerated search
spaces, no one-vertex witness appears and the two-vertex pattern is the first
observed regime where the separation occurs.  That latter statement remains
certificate-level evidence, not a proof.

\subsection{Odd-$k$ Separation (\Cref{thm:odd-k-separation})}

The proof in the main text is self-contained.  Here we provide additional
commentary on the boundary cases.

\paragraph{The case $m = 0$ ($k = 1$).}
When $m = 0$, the sample $S_0 = ((a,0), (b,0))$ reduces to the two-point
witness of \Cref{sec:minimal-1nn}.  The margin analysis gives
$\M_1(S_0,a) = -1$; after LOO at index $0$ the margin is $-1$, and after
replacement at index $0$ the margin is $+1$.  This $k=1$ case is excluded from
\Cref{thm:odd-k-separation} (which requires $m \ge 1$, i.e., $k \ge 3$) and is
handled directly by the two-point witness.

\paragraph{The case $m = 1$ ($k = 3$).}
The sample is $S_1 = ((a,0), (a,0), (a,1), (b,0))$ of size $n=4$.  The original
margin at $a$ is $-1$.  LOO at each index is $0$ by the case analysis in the
proof.  After replacing the first $(a,0)$ by $(a,1)$, the margin becomes $+1$.
This matches the computational candidate found by the earlier bounded search.

\paragraph{General $m \ge 2$.}
For $m \ge 2$, $n = 3m+1 \ge 7$ and the post-deletion sample size $3m$ is
always at least $k = 2m+1$, so $k' = k$ throughout.  The LOO analysis is
uniform across all $m$.

\subsection{Even-$k$ Separation (\Cref{thm:even-k-separation})}

The even-$k$ case follows the same structure but with margin $0$ in the
original configuration.  The tie-breaking rule $0 \prec 1$ determines the
prediction in the pre- and post-deletion configurations.  If the tie-breaking
rule were changed to $1 \prec 0$, the separation would reverse direction.
